% paper-skeleton.tex — minimal acmart submission skeleton for the Week 14 Draft Paper.
% Swap the \documentclass line for IEEEtran if your target venue is IEEE:
%   \documentclass[conference]{IEEEtran}
% acmart.cls v2.18 (2026/06/01) and IEEEtran.cls v1.8b are pre-loaded on Overleaf.
% Single-column "manuscript,review" is the ACM submission format; "anonymous" -> double-blind.

\documentclass[manuscript,review,anonymous]{acmart}

\begin{document}

\title{<Your title: a claim, not a topic>}
% Authors are hidden by "anonymous" for double-blind review; fill at camera-ready.

\begin{abstract}
% FIVE MOVES (chapter 14.2). Put a NUMBER in the result move.
% [CONTEXT]      The world the problem lives in.
% [GAP]          What is missing/wrong (the cited absence).
% [APPROACH]     What you did that is new (the startling sentence).
% [RESULT]       The finding, with effect size / CI.
% [SIGNIFICANCE] Why a reader should care.
\end{abstract}

\maketitle

\section{Introduction}
% CARS, three moves (14.3): territory -> niche/gap -> occupy. End with a contributions list:
% This paper makes the following contributions:
%   C1. ...
%   C2. ...
%   C3. ...

\section{Related Work}
% POSITION, don't summarize (14.4). Group by approach; pull from your comparison matrix.
% "Retrieval-based methods [x,y] achieve A but assume B; in contrast, we ..."

\section{Methodology}
% WHAT YOU DID, reproducibly (14.5). Data + splits, baselines, ablations, metric (+ construct),
% seeds/runs, hardware, software versions. If an LLM was an instrument: model NAME, VERSION,
% DATE, decoding settings (e.g. reasoning effort), prompt (long prompts -> appendix). No spin, no results here.

\section{Results}
% WHAT YOU FOUND (14.5). Tables + figures carry the load. Effect sizes + CIs, not bare p-values.
% Report the spread across seeds. EVERY run goes in. No interpretation in this section.

\section{Discussion}
% WHAT IT MEANS (14.5, ch.13). Interpret, relate to prior work, limitations + threats to validity
% (Cook & Campbell's four categories). Do NOT overclaim past what the CIs/effect sizes support.

\section{Conclusion and Future Work}
% The takeaway and the open doors. State the contribution plainly; don't hedge it away.

\bibliographystyle{ACM-Reference-Format}   % use IEEEtran for IEEE
\bibliography{references}                   % every entry verified against dblp / Semantic Scholar

\end{document}
